% =============================================================
%  Tugas Individu 02 — Representasi Teks
%  IF25-40304 · Pembelajaran Mesin Multimodal · ITERA
% =============================================================
\documentclass[12pt, a4paper]{article}

\usepackage{tugas-style}

\renewcommand{\assignmentnumber}{02}
\renewcommand{\doctitle}{Tugas Individu 02}
\renewcommand{\duedate}{\color{ctpBlue}\faCalendarCheck\ \ Tenggat: sesuai jadwal kelas \faCalendarCheck}

\begin{document}

% ── HEADER ──────────────────────────────────────────────────
\makeassignmentheader

\hrule width0.3\textwidth
\revisioninfo{%
  \vhEntry{0.1}{16 September 2026}{I.W.W.}{Draf inisial tugas.}
}
\hrule
\vskip .3in

\tableofcontents
\clearpage

% =============================================================
\section{Tujuan Pembelajaran}\label{sec:objectives}

Setelah menyelesaikan tugas ini, mahasiswa mampu:
\begin{objectives}
  \item Membedakan representasi \emph{one-hot} (sparse, ortogonal) dan \emph{word embeddings} (dense, semantik).
  \item Menghitung kemiripan kosinus dan menginterpretasikan kedekatan vektor kata.
  \item Menggunakan operasi aritmetika vektor untuk menguji relasi semantik (analogi).
  \item Memproyeksikan ruang vektor kata berdimensi tinggi ke 2D menggunakan PCA.
  \item Menjelaskan keterbatasan embeddings statis terhadap polisemi dan keunggulan contextual embeddings.
\end{objectives}

% =============================================================
\section{Catatan dan Batasan}\label{sec:notes}
\vspace{-.1cm}

\begin{minipage}{\linewidth}
\begin{bclogo}[couleur=ctpBase, arrondi=0.1, logo=\bccrayon, ombre=false, couleurBord=ctpSurface, epBord=0.6]{Perhatikan hal berikut:}
  \begin{coloredPen}
    \item Tugas ini bersifat \textbf{individu}; tidak diperkenankan bekerja sama atau menyalin jawaban.
    \item Kerjakan eksperimen secara \textbf{inferensi saja} --- tidak ada pelatihan model baru.
    \item Tulis analisis dengan bahasa sendiri, bukan sekadar menyalin output program.
    \item Model fastText berukuran besar ($\pm$1\,GB); unduh lebih awal atau gunakan alternatif yang disediakan.
  \end{coloredPen}
\end{bclogo}
\end{minipage}

% =============================================================
\section{Perangkat Lunak dan Peralatan yang Diperlukan}\label{sec:requiredsw}

\begin{itemize}[itemsep=2pt,parsep=0pt,topsep=2pt,partopsep=2pt]
  \item[\color{ctpBlue}\faPython] \textbf{Python} 3.8 atau lebih baru.
  \item[\color{ctpBlue}\faFileCode] \textbf{Paket wajib:} \code{numpy}, \code{matplotlib}, \code{scikit-learn}, \code{gensim}.
  \item[\color{ctpBlue}\faFileCode] \textbf{Paket opsional (Soal 5 bonus):} \code{transformers}, \code{torch}.
  \item[\color{ctpBlue}\faDownload] \textbf{Model pra-latih:} fastText bahasa Indonesia (\code{cc.id.300.vec.gz}, $\pm$1\,GB).
\end{itemize}

\begin{codewindow}{setup.sh}
\begin{lstlisting}[style=pycode]
pip install numpy matplotlib scikit-learn gensim
\end{lstlisting}
\end{codewindow}

\begin{minipage}{\linewidth}
\begin{bclogo}[couleur=ctpMantle, arrondi=0.1, logo=\bcinfo, ombre=false, couleurBord=ctpSurface, epBord=0.6]{Catatan Model Pra-Latih}
  Model fastText Indonesia (\code{cc.id.300.vec.gz}) berukuran $\pm$1\,GB. Unduh dari
  \href{https://fasttext.cc/docs/en/crawl-vectors.html}{fasttext.cc}. Jika terkendala
  bandwidth, gunakan \code{glove-wiki-gigaword-50} via \code{gensim.downloader} sebagai
  alternatif ($\pm$66\,MB, kosakata bahasa Inggris).
\end{bclogo}
\end{minipage}

% =============================================================
\section{Pernyataan Masalah}\label{sec:intro}

Tugas ini melengkapi materi Pertemuan 02 tentang representasi teks untuk \emph{deep learning}.
Anda akan melakukan eksperimen komputasi dan analisis tertulis untuk memverifikasi tiga klaim yang dibahas pada slide kuliah:
\begin{enumerate}
  \item One-hot encoding bersifat \textbf{ortogonal} --- tidak menangkap kemiripan makna antarkata.
  \item Word embeddings bersifat \textbf{dense dan semantik} --- kata bermakna mirip berdekatan dalam ruang vektor.
  \item Embeddings statis memiliki limitasi terhadap \textbf{polisemi} --- satu kata yang sama selalu mendapat vektor identik, terlepas dari konteks.
\end{enumerate}

Rincian dan persyaratan pengerjaan diberikan pada bagian berikutnya.

% =============================================================
\section{Persyaratan}\label{sec:requirements}

Perhatikan setiap persyaratan pada soal-soal berikut.

% ------------------------------------------------------------
\subsection{Soal 1 --- One-Hot Encoding \& Keterbatasannya}
\label{sec:soal1}

Representasi teks paling sederhana adalah one-hot encoding, di mana setiap kata direpresentasikan sebagai vektor biner berdimensi = ukuran kosakata. Vektor ini bersifat \emph{ortogonal}: setiap pasang kata berbeda menghasilkan kemiripan kosinus nol.

\subsubsection{Instruksi}
\label{sec:soal1-instruksi}

\textbf{(a)} Diberikan kosakata berikut:

\begin{codewindow}{soal1\_data.py}
\begin{lstlisting}[style=pycode]
vocab = ["presiden", "teknologi", "musik", "indonesia", "sains"]
\end{lstlisting}
\end{codewindow}

Bangun fungsi \code{one\_hot(kata, vocab)} yang mengembalikan vektor one-hot untuk sebuah kata. Tampilkan vektor untuk setiap kata dalam kosakata.

\textbf{(b)} Implementasikan fungsi cosine similarity secara manual:

\begin{center}
\colorbox{ctpBase}{%
  \parbox{0.75\linewidth}{\centering
    \footnotesize\bfseries\color{ctpOverlay}\MakeUppercase{Cosine Similarity}\\[6pt]
    {\large\color{ctpBlue}$\displaystyle\operatorname{sim}(\mathbf{a}, \mathbf{b}) = \frac{\mathbf{a} \cdot \mathbf{b}}{\|\mathbf{a}\| \cdot \|\mathbf{b}\|}$}}}
\end{center}

Hitung dan tampilkan \textbf{matriks kemiripan kosinus ($5 \times 5$)} antar semua pasang kata.

\textbf{(c)} Jawab secara tertulis (3--5 kalimat):
\begin{itemize}
  \item Mengapa semua kemiripan antar kata berbeda bernilai \textbf{0.0}?
  \item Apa implikasinya terhadap model ML yang menerima input one-hot?
  \item Secara intuitif, pasangan kata mana yang \emph{seharusnya} lebih mirip secara semantik? Mengapa one-hot tidak bisa menangkap ini?
\end{itemize}

\subsubsection{Kerangka Kode}
\label{sec:soal1-kode}

\begin{codewindow}{soal1\_kerangka.py}
\begin{lstlisting}[style=pycode]
import numpy as np

vocab = ["presiden", "teknologi", "musik", "indonesia", "sains"]
word_to_id = {kata: i for i, kata in enumerate(vocab)}

def one_hot(kata, vocab, word_to_id):
    """Mengembalikan vektor one-hot untuk sebuah kata."""
    vektor = np.zeros(len(vocab))
    vektor[word_to_id[kata]] = 1
    return vektor

def cosine_similarity(a, b):
    """Kemiripan kosinus: (a.b) / (||a|| . ||b||)."""
    return np.dot(a, b) / (np.linalg.norm(a) * np.linalg.norm(b))
\end{lstlisting}
\end{codewindow}

% ------------------------------------------------------------
\clearpage
\subsection{Soal 2 --- Word Embeddings: Vektor Dense \& Kemiripan Semantik}
\label{sec:soal2}

Word embeddings menggantikan vektor one-hot yang panjang dan \emph{sparse} dengan vektor \emph{dense} berdimensi rendah yang merepresentasikan makna. Kata yang bermakna mirip akan saling berdekatan dalam ruang vektor.

\subsubsection{Instruksi}
\label{sec:soal2-instruksi}

\textbf{(a)} Muat model pra-latih fastText bahasa Indonesia:

\begin{codewindow}{soal2\_muat.py}
\begin{lstlisting}[style=pycode]
from gensim.models import KeyedVectors

model = KeyedVectors.load_word2vec_format("cc.id.300.vec.gz", binary=False)
\end{lstlisting}
\end{codewindow}

Untuk setiap kata berikut, tampilkan \textbf{5 kata paling mirip} beserta skor kemiripan kosinus: \code{"presiden"}, \code{"teknologi"}, \code{"musik"}.

\begin{codewindow}{soal2\_similar.py}
\begin{lstlisting}[style=pycode]
for kata in ["presiden", "teknologi", "musik"]:
    print(f"\n5 kata paling mirip dengan '{kata}':")
    for kata_mirip, skor in model.most_similar(kata, topn=5):
        print(f"  {kata_mirip:<16} {skor:.4f}")
\end{lstlisting}
\end{codewindow}

\textbf{(b)} Hitung dan tampilkan secara numerik:

\begin{codewindow}{soal2\_numerik.py}
\begin{lstlisting}[style=pycode]
sim_presiden_indonesia = model.similarity("presiden", "indonesia")
sim_presiden_musik     = model.similarity("presiden", "musik")

print(f"sim(presiden, indonesia) = {sim_presiden_indonesia:.4f}")
print(f"sim(presiden, musik)     = {sim_presiden_musik:.4f}")
\end{lstlisting}
\end{codewindow}

\textbf{(c)} Jawab dalam 3--5 kalimat:
\begin{itemize}
  \item Apakah \code{"presiden"} lebih dekat ke \code{"indonesia"} atau ke \code{"musik"} di ruang embedding?
  \item Bandingkan secara eksplisit dengan hasil Soal 1: apa yang berubah ketika berpindah dari one-hot ke embeddings?
  \item Bagaimana embeddings mengatasi keterbatasan ortogonalitas?
\end{itemize}

% ------------------------------------------------------------
\subsection{Soal 3 --- Analogi Vektor dalam Bahasa Indonesia}
\label{sec:soal3}

Sifat penting word embeddings adalah kemampuannya merepresentasikan \emph{relasi semantik} melalui operasi aritmetika vektor. Operasi seperti \code{raja - pria + wanita = ratu} menunjukkan bahwa embeddings menangkap hubungan gender, geografi, dan peran secara geometris.

\subsubsection{Instruksi}
\label{sec:soal3-instruksi}

Uji analogi vektor berikut pada fastText Indonesia:

\begin{center}\small
\begin{tabular}{@{} C{0.6cm} L{6.5cm} L{4cm} @{}}
  \toprule \textbf{No} & \textbf{Rumus Analogi} & \textbf{Hasil Diharapkan} \\
  \midrule
  (a) & \code{raja - pria + wanita = ?}          & ratu \\
  (b) & \code{jepang - asia + eropa = ?}         & jerman / perancis \\
  (c) & \code{jakarta - indonesia + jepang = ?}  & tokyo \\
  \bottomrule
\end{tabular}
\end{center}

\begin{codewindow}{soal3\_analogi.py}
\begin{lstlisting}[style=pycode]
# (a) raja - pria + wanita
hasil_a = model.most_similar(
    positive=["wanita", "raja"], negative=["pria"], topn=3
)
print("raja - pria + wanita =")
for kata, skor in hasil_a:
    print(f"  {kata:<16} {skor:.4f}")

# (b) jepang - asia + eropa
hasil_b = model.most_similar(
    positive=["eropa", "jepang"], negative=["asia"], topn=3
)
print("\njepang - asia + eropa =")
for kata, skor in hasil_b:
    print(f"  {kata:<16} {skor:.4f}")

# (c) jakarta - indonesia + jepang
hasil_c = model.most_similar(
    positive=["jepang", "jakarta"], negative=["indonesia"], topn=3
)
print("\njakarta - indonesia + jepang =")
for kata, skor in hasil_c:
    print(f"  {kata:<16} {skor:.4f}")
\end{lstlisting}
\end{codewindow}

Tampilkan hasil \textbf{top-3} untuk setiap analogi, lalu jawab:
\begin{itemize}
  \item Apakah ketiga analogi menghasilkan kata yang masuk akal?
  \item Jika hasilnya tidak sempurna, hipotesiskan mengapa (misal: bias korpus, frekuensi kata, ambiguitas).
\end{itemize}

\begin{callout}[ctpGreen]
\textbf{\color{ctpGreen}Bonus:} Buat \textbf{satu analogi kreatif} Anda sendiri dalam bahasa Indonesia. Pilih kata-kata yang memiliki hubungan semantik jelas. Tampilkan hasil dan evaluasi keberhasilannya dalam 2--3 kalimat.
\end{callout}

% ------------------------------------------------------------
\clearpage
\subsection{Soal 4 --- Visualisasi Ruang Vektor Kata Indonesia}
\label{sec:soal4}

Untuk mengamati sifat ``kata bermakna mirip berdekatan'', vektor embedding diproyeksikan ke dua dimensi menggunakan PCA (\emph{Principal Component Analysis}).

\subsubsection{Instruksi}
\label{sec:soal4-instruksi}

Pilih \textbf{12 kata} dari fastText Indonesia yang terdiri dari \textbf{3--4 kelompok semantik}:

\begin{center}\small
\begin{tabular}{@{} L{2.5cm} L{8cm} @{}}
  \toprule \textbf{Kelompok} & \textbf{Contoh Kata} \\
  \midrule
  Negara  & \code{indonesia}, \code{jepang}, \code{brasil} \\
  Hewan   & \code{kucing}, \code{anjing}, \code{burung} \\
  Profesi & \code{dokter}, \code{guru}, \code{insinyur} \\
  Makanan & \code{nasi}, \code{rendang}, \code{sate} \\
  \bottomrule
\end{tabular}
\end{center}

\textbf{(a)} Ekstrak vektor dan reduksi dimensi ke 2D menggunakan PCA, lalu buat scatter plot:

\begin{codewindow}{soal4\_visualisasi.py}
\begin{lstlisting}[style=pycode]
import numpy as np
from sklearn.decomposition import PCA
import matplotlib.pyplot as plt

kata_pilihan = [
    "indonesia", "jepang", "brasil",
    "kucing", "anjing", "burung",
    "dokter", "guru", "insinyur",
    "nasi", "rendang", "sate",
]

kelompok = {
    "indonesia": "Negara",   "jepang": "Negara",   "brasil": "Negara",
    "kucing": "Hewan",       "anjing": "Hewan",     "burung": "Hewan",
    "dokter": "Profesi",     "guru": "Profesi",     "insinyur": "Profesi",
    "nasi": "Makanan",       "rendang": "Makanan",  "sate": "Makanan",
}

warna = {
    "Negara": "#1e66f5", "Hewan": "#40a02b",
    "Profesi": "#8839ef", "Makanan": "#fe640b",
}

# Ekstrak vektor & reduksi dimensi
vektor = np.array([model[k] for k in kata_pilihan])
pca = PCA(n_components=2)
koordinat_2d = pca.fit_transform(vektor)

# Plot
plt.figure(figsize=(8, 6))
sudah_label = set()
for i, kata in enumerate(kata_pilihan):
    k = kelompok[kata]
    lbl = k if k not in sudah_label else ""
    sudah_label.add(k)
    plt.scatter(koordinat_2d[i, 0], koordinat_2d[i, 1],
                c=warna[k], s=100, label=lbl)
    plt.annotate(kata, (koordinat_2d[i, 0], koordinat_2d[i, 1]),
                 textcoords="offset points", xytext=(8, 4),
                 fontsize=10, fontweight="bold")
plt.title("Proyeksi 2D Vektor Kata Indonesia (PCA)")
plt.xlabel("Komponen Utama 1")
plt.ylabel("Komponen Utama 2")
plt.legend()
plt.grid(alpha=0.3)
plt.tight_layout()
plt.savefig("visualisasi_embeddings.png", dpi=150)
plt.show()
\end{lstlisting}
\end{codewindow}

\textbf{(b)} Analisis (3--5 kalimat):
\begin{itemize}
  \item Apakah kluster yang terbentuk sesuai intuisi Anda?
  \item Kata mana yang paling ``tercampur'' atau berjauhan dari klompoknya?
  \item Mengapa menurut Anda hal tersebut terjadi?
\end{itemize}

% ------------------------------------------------------------
\subsection{Soal 5 --- Polisemi: Limitasi Embeddings Statis}
\label{sec:soal5}

Word embeddings klasik memberikan \textbf{satu vektor tetap per kata}, sehingga tidak dapat membedakan makna kata yang bergantung pada konteks. Kata \textbf{``tahu''} dalam bahasa Indonesia memiliki dua makna yang benar-benar berbeda: makanan dari kedelai (nomina) dan mengetahui (verba).

\subsubsection{Instruksi}
\label{sec:soal5-instruksi}

\textbf{(a)} Dari fastText Indonesia, ambil vektor kata \code{"tahu"}. Tampilkan 5 kata terdekat:

\begin{codewindow}{soal5\_polisemi.py}
\begin{lstlisting}[style=pycode]
print("5 kata paling mirip dengan 'tahu':")
for kata, skor in model.most_similar("tahu", topn=5):
    print(f"  {kata:<16} {skor:.4f}")
\end{lstlisting}
\end{codewindow}

\textbf{(b)} Kata \code{"tahu"} memiliki \textbf{2 makna yang berbeda}:

\begin{callout}[ctpBlue]
{\color{ctpBlue}\bfseries Makanan:} ``Ibu membeli \textbf{tahu} goreng di pasar.'' (makanan dari kedelai)
\end{callout}
\begin{callout}[ctpBlue]
{\color{ctpBlue}\bfseries Mengetahui:} ``Saya \textbf{tahu} jawaban soal itu.'' (kata kerja: mengetahui)
\end{callout}

Jawab dalam 3--5 kalimat:
\begin{itemize}
  \item Apakah 5 kata terdekat dari fastText mencerminkan \textbf{kedua} makna tersebut, atau hanya salah satu?
  \item Apakah makna yang terwakili cenderung ke makanan atau ke pengetahuan? Mengapa?
  \item Mengapa hal ini menjadi keterbatasan fundamental embeddings statis?
\end{itemize}

\textbf{(c)} \emph{(Opsional --- Bonus)} Jika lingkungan mendukung \code{transformers} dan \code{torch}:

\begin{codewindow}{soal5\_bonus\_indobert.py}
\begin{lstlisting}[style=pycode]
from transformers import AutoTokenizer, AutoModel
import torch

# Muat IndoBERT pra-latih (inferensi saja, tanpa pelatihan)
name = "indobenchmark/indobert-base-p1"
tokenizer = AutoTokenizer.from_pretrained(name)
model_bert = AutoModel.from_pretrained(name)
model_bert.eval()

def vektor_kata(kalimat, kata):
    """Vektor kontekstual dari hidden state BERT."""
    enc = tokenizer(kalimat, return_tensors="pt")
    with torch.no_grad():
        out = model_bert(**enc)
    hidden = out.last_hidden_state[0]
    ids = tokenizer.encode(kata, add_special_tokens=False)
    pos = (enc["input_ids"][0] == ids[0]).nonzero()[0].item()
    return hidden[pos]

# Kata yang sama, dua konteks berbeda
s1 = "Ibu membeli tahu goreng di pasar."
s2 = "Saya tahu jawaban soal itu."

v1 = vektor_kata(s1, "tahu")
v2 = vektor_kata(s2, "tahu")

cos = torch.nn.functional.cosine_similarity(
    v1.unsqueeze(0), v2.unsqueeze(0)
).item()

print(f"Kemiripan vektor 'tahu' antar dua konteks: {cos:.4f}")
print("Nilai jauh dari 1.0 membuktikan embeddings kontekstual berbeda.")
\end{lstlisting}
\end{codewindow}

Bandingkan hasil IndoBERT dengan fastText. Tulis kesimpulan:
\begin{itemize}
  \item Seberapa berbeda vektor kontekstual \code{"tahu"} di dua kalimat tersebut?
  \item Bagaimana hal ini membuktikan keunggulan contextual embeddings?
\end{itemize}

% ------------------------------------------------------------
\clearpage
\subsection{Soal 6 --- Refleksi: Evolusi Arsitektur Sekuensial}
\label{sec:soal6}

Representasi teks yang baik memerlukan arsitektur yang mampu menangkap \textbf{ketergantungan sekuensial}. Bagian ini menguji pemahaman konseptual terhadap evolusi arsitektur dari RNN ke Transformer.

\subsubsection{Instruksi}
\label{sec:soal6-instruksi}

Jawab \textbf{dua} pertanyaan berikut dalam bentuk paragraf singkat (masing-masing \textbf{5--8 kalimat}):

\begin{callout}[ctpLavender]
\textbf{\color{ctpLavender}(a) Vanishing Gradient \& Mekanisme LSTM}\par\vspace{2pt}
RNN memiliki masalah \emph{vanishing gradient}: saat memproses sekuens panjang, informasi dari kata-kata awal menghilang sebelum mencapai kata akhir. Jelaskan secara intuitif:
\begin{itemize}
  \item Mengapa masalah ini muncul? (Kaitkan dengan propagasi gradien melalui banyak langkah waktu.)
  \item Bagaimana LSTM mengatasinya melalui mekanisme \emph{gates} (forget gate, input gate, output gate)?
  \item Gunakan \textbf{analogi sederhana} untuk memperjelas penjelasan Anda.
\end{itemize}
\end{callout}

\begin{callout}[ctpLavender]
\textbf{\color{ctpLavender}(b) RNN/LSTM vs Transformer: Trade-off}\par\vspace{2pt}
Transformer kini mendominasi NLP melalui mekanisme \emph{self-attention} yang memproses seluruh sekuens secara paralel. Namun, RNN/LSTM masih relevan di beberapa skenario.
\begin{itemize}
  \item Sebutkan \textbf{satu contoh konkret} skenario di mana RNN/LSTM mungkin lebih unggul atau lebih praktis daripada Transformer.
  \item Jelaskan alasannya dengan merujuk pada \textbf{sifat arsitektur} keduanya.
\end{itemize}
\end{callout}

% =============================================================
\clearpage
\section{Kriteria Penilaian}\label{sec:evaluation}

Tugas Anda akan dinilai berdasarkan kriteria berikut:

\renewcommand{\arraystretch}{1.5}
\begin{center}\small
\begin{tabular}{|p{11.5cm}|c|}
  \hline
  \thead{\color{ctpBlue} Kriteria} & \thead{\color{ctpBlue}Bobot} \\
  \hline
  Kebenaran teknis --- kode berjalan tanpa error, output benar, penggunaan library tepat. & 40\% \\
  \hline
  Kedalaman analisis --- penjelasan tertulis mendalam, bukan sekadar mendeskripsikan output. & 30\% \\
  \hline
  Eksplorasi kreatif --- pemilihan kata/kelompok menarik (Soal 4), analogi kreatif (Soal 3). & 15\% \\
  \hline
  Kualitas presentasi --- laporan rapi, terstruktur, visualisasi informatif dan berlabel. & 15\% \\
  \hline
  \textbf{Total} & \textbf{100\%} \\
  \hline
\end{tabular}
\end{center}

% =============================================================
\section{Apa yang Dikumpulkan}\label{sec:submit}

Seluruh artefak tugas dikemas ke dalam \textbf{satu berkas arsip} dengan format penamaan:
\begin{center}
  \code{mgg02\_tugas\_nim.zip}
\end{center}

dengan \code{nim} diganti oleh Nomor Induk Mahasiswa (NIM) masing-masing tanpa spasi.

\textbf{Artefak yang dapat disertakan di dalam arsip} (dapat lebih dari satu jenis berkas, disesuaikan dengan bentuk penyelesaian Anda):
\begin{borderedsquare}
  \item \code{mgg02\_tugas\_NIM.ipynb} --- apabila penyelesaian dikerjakan menggunakan Jupyter Notebook (termasuk Google Colab). Notebook \textbf{wajib telah dijalankan} sehingga seluruh \emph{output} sel tersimpan di dalam berkas.
  \item \code{mgg02\_tugas\_NIM.py} --- apabila penyelesaian dikerjakan dalam bentuk skrip Python.
  \item \code{mgg02\_tugas\_NIM.pdf} --- apabila analisis ditulis dalam dokumen teks terpisah.
  \item \code{mgg02\_tugas\_xx.jpg} --- apabila terdapat plot atau gambar pendukung tambahan; \code{xx} merupakan nomor urut visualisasi (misalnya \code{01}, \code{02}, dan seterusnya).
\end{borderedsquare}

Setiap berkas menggunakan awalan \code{mgg02\_tugas\_NIM} yang sama, dengan \code{NIM} ditulis secara konsisten (huruf besar/kecil mengikuti NIM resmi). Pastikan seluruh berkas berada tepat di dalam satu arsip \code{.zip} tanpa struktur folder bertingkat yang tidak perlu.

\begin{callout}[ctpBlue]
\textbf{\color{ctpBlue}Catatan Pengerjaan}
\begin{itemize}
  \item \textbf{Soal 1--5:} Wajib dikerjakan oleh semua mahasiswa.
  \item \textbf{Soal 5(c):} Opsional / Bonus --- memerlukan \code{transformers} + \code{torch} ($\sim$400\,MB).
  \item \textbf{Soal 6:} Dapat dikerjakan tanpa kode --- murni pemahaman konseptual.
  \item \textbf{Estimasi waktu:} 3--4 jam.
\end{itemize}
\end{callout}

\section*{Referensi}
{\small\color{ctpSubtext}
\begin{enumerate}[label={[\arabic*]}, leftmargin=2em]
  \item Mikolov, T., et al.\ (2013). \emph{Efficient Estimation of Word Representations in Vector Space}.
  \item Vaswani, A., et al.\ (2017). \emph{Attention Is All You Need}.
  \item Devlin, J., et al.\ (2019). \emph{BERT: Pre-training of Deep Bidirectional Transformers}.
  \item Bojanowski, P., et al.\ (2017). \emph{Enriching Word Vectors with Subword Information} (fastText).
  \item Wilie, B., et al.\ (2020). \emph{IndoNLU: Benchmark and Resources for Evaluating Indonesian NLU}.
\end{enumerate}}

\end{document}
